\section{Method}
\label{sec:method}

\subsection{Data sources}
\label{sec:data}

The training corpus is sourced from four primary works, summarized in
Table~\ref{tab:sources}. Two are English public-domain translations from
the early twentieth century, one is a contemporary copyrighted catechetical
work used here under a fair-use research posture, and one (City of God)
shares the same translator family as Confessions.

\begin{table}[H]
\centering
\caption{Primary sources used to build the training corpus.}
\label{tab:sources}
\small
\setlength{\tabcolsep}{6pt}
\renewcommand{\arraystretch}{1.2}
\begin{tabular}{p{4.4cm} p{4.0cm} p{3.2cm} p{2.0cm}}
\toprule
\textbf{Work} & \textbf{Translation} & \textbf{Status} & \textbf{Source} \\
\midrule
Catechism of the Catholic Church (1992) & USCCB / Libreria Editrice Vaticana & Copyrighted (fair use) & vatican.va \\
\emph{Summa Theologica} & Fr.~Laurence \mbox{Shapcote}, 1920 & Public domain (US) & newadvent.org \\
Augustine, \emph{Confessions} & E.B.~Pusey & Public domain & newadvent.org \\
Augustine, \emph{City of God} & Marcus Dods & Public domain & newadvent.org \\
\bottomrule
\end{tabular}
\end{table}

After scraping with a respectful crawler (descriptive user agent,
1-second sleep between requests, on-disk cache, \texttt{robots.txt}
compliance) and cleaning, we obtain a corpus of 342 structured chunks
distributed across the four sources. Neither the raw HTML, the cleaned
JSONL chunks, nor the trained adapter weights are committed to the public
repository; only code, recipe, aggregate metrics, and short qualitative
excerpts are shared. See the project's \texttt{DATA\_LICENSING.md} for
details.

\subsection{Pipeline}
\label{sec:pipeline}

Figure~\ref{fig:pipeline} shows the end-to-end data path. The pipeline is
deliberately small and inspectable, each arrow is a short Python script
with a single responsibility.

\begin{figure}[H]
\centering
\begin{tikzpicture}[
    >=Stealth,
    box/.style={
      rectangle,
      rounded corners=3pt,
      draw=black!60,
      fill=blue!4,
      minimum height=1.45cm,
      minimum width=2.7cm,
      align=center,
      inner sep=4pt,
      text width=2.5cm,
      font=\small
    },
    arrow/.style={->, thick, draw=black!60, shorten <=1pt, shorten >=1pt},
    node distance=0.4cm
  ]

  \node[box] (sources) {\textbf{Primary sources}\\[2pt]\footnotesize CCC, \emph{Summa},\\ Augustine};
  \node[box, right=of sources] (scrape)  {\textbf{Scrape}\\[2pt]\footnotesize Rate-limited\\ HTTP + cache};
  \node[box, right=of scrape]  (clean)   {\textbf{Clean}\\[2pt]\footnotesize HTML\,$\to$\,JSONL\\ chunks};
  \node[box, right=of clean]   (gen)     {\textbf{Distill}\\[2pt]\footnotesize Claude Sonnet\\ 4.6 teacher};
  \node[box, right=of gen]     (sft)     {\textbf{SFT}\\[2pt]\footnotesize Qwen 2.5\,7B\\ LoRA, MLX Q8};

  \draw[arrow] (sources) -- (scrape);
  \draw[arrow] (scrape)  -- (clean);
  \draw[arrow] (clean)   -- (gen);
  \draw[arrow] (gen)     -- (sft);

\end{tikzpicture}

\caption{End-to-end data pipeline for Phase~1. Primary sources are scraped
respectfully, cleaned into a structured JSONL chunk corpus, then expanded
by a Claude Sonnet 4.6 teacher into (question, scholastic-answer) pairs
that drive LoRA SFT of Qwen 2.5\,7B-Instruct under MLX with 8-bit weight
quantization.}
\label{fig:pipeline}
\end{figure}

\subsection{Teacher-distilled training data}
\label{sec:teacher}

For each source chunk we ask Claude Sonnet 4.6, acting as a teacher model,
to produce 2--3 (question, answer) pairs in which the answer is written
in scholastic register. The system prompt (paraphrased here; full text in
\texttt{scripts/generate\_training\_pairs.py}) instructs the teacher to:

\begin{itemize}
  \item produce a question a graduate student of theology or philosophy
        might plausibly ask, drawn from the chunk's subject matter;
  \item answer in either Summa form (``It seems that\,\dots Objection
        1\,\dots On the contrary\,\dots I answer that\,\dots Reply to
        Objection 1\,\dots'') or in Augustinian form (autobiographical
        second-person address with scriptural allusion), as fits the
        chunk;
  \item cite the CCC by paragraph number where doctrinally apposite;
  \item refuse to invent CCC paragraph numbers if not directly attested in
        the chunk, instead citing thematically.
\end{itemize}

The teacher outputs JSON conforming to a strict schema, validated by the
client. Across 342 chunks, the teacher produced 93 (question, answer)
pairs (we discard malformed JSON or pairs shorter than a threshold). We
split 83/10 train/valid by random shuffle with a fixed seed. The
total API cost for this run was approximately \$0.50 at then-current
Claude Sonnet pricing.

\subsection{Model and training}
\label{sec:training}

We start from Qwen~2.5\,7B-Instruct \cite{qwen25}, converted to MLX format
with 8-bit weight quantization
(\texttt{mlx\_lm.convert --q --q-bits 8}). On disk the quantized model is
$\approx$\,7.8\,GB; resident memory during training peaks at 12.7\,GB.

We attach LoRA adapters to the top 16 transformer layers (out of 28),
training with AdamW at a learning rate of $10^{-5}$, batch size~1, maximum
sequence length~2048, for 200 iterations. With batch~1 and 83 training
pairs this is approximately 2.4 epochs. Training throughput is
0.48~it/s ($\approx$~235~tok/s) and the full 200-iteration run takes about
seven minutes wall-clock on the M4~Pro. We use the standard MLX-LM LoRA
\cite{mlx} training entrypoint (\texttt{mlx\_lm\_lora.train}). The
adapter writes 2.6M trainable parameters out of 7.6B total
(0.034\,\%).

\subsection{Evaluation rubric}
\label{sec:rubric}

Because no shared benchmark exists for scholastic-register transfer, we
score model outputs on a custom rubric implemented in
\texttt{src/scholastic/rubric.py}. The rubric is intentionally simple:
four orthogonal dimensions, each scored 0--3, giving a per-prompt total
0--12. Across $N$ prompts the maximum is $12N$.

\begin{description}
  \item[Scholastic register (0--3)] Counts of distinct surface markers
        like \emph{I answer that}, \emph{It seems that}, \emph{Objection},
        \emph{Reply to Objection}, \emph{On the contrary}, and
        scholastic connectives (\emph{accordingly}, \emph{insofar as},
        \emph{wherein}, etc.).
  \item[Augustinian voice (0--3)] Distinct hits among autobiographical or
        confessional markers (\emph{O Lord}, archaic second-person
        pronouns, \emph{my soul}, \emph{the City of God}).
  \item[CCC grounding (0--3)] Visible paragraph-number citations (\S~$n$,
        \emph{CCC $n$}, ``paragraph $n$'').
  \item[Structure (0--3)] Number of paragraphs combined with the
        presence of objection-and-response structural markers.
\end{description}

This rubric is a coarse signal. It rewards lexical and structural
surface form rather than theological correctness, and it can be gamed by
shallow string interpolation. We treat it as a useful proxy for
register transfer, not as a substitute for human evaluation; see
Section~\ref{sec:limitations} for discussion.
